\section{Introduction}
The success of modern neural networks (NNs) is often attributed to feature learning, namely the ability of models to adapt to the structure of the data during training \citep{bach2017breaking,suzuki2019adaptivity,damian2022neural}. This stands in contrast to non-adaptive approaches, such as kernel methods and neural networks operating in the so-called lazy training regime \citep{jacot2018neural,chizat2019lazy}, where features remain effectively frozen at their random, data-independent initialization.

\if 0
From the perspective of function space, we consider function classes generated by nonlinear feature maps. Let $\mathcal{X}$ and $\mathcal{W}$ be compact sets, $\rho$ and $\tau$ be probability measures on $\mathcal{X}$ and $\mathcal{W}$, respectively, and $\phi: \mathcal{X} \times \mathcal{W} \rightarrow \mathbb{R}$ be a bounded feature map.
We define the operator $T : L^2(\mathrm{d}\tau) \to L^2(\mathrm{d}\rho)$ as
\begin{equation}\label{opetatorT}
(Ta)(\bm x) = \int_{\mathcal{W}} a(\bm w)\phi(\bm x, \bm w)\, \mathrm{d} \tau(\bm w)\,.
\end{equation}
It yields several examples. A classical one is the reproducing kernel Hilbert space (RKHS) $\mathcal{F}_2$, obtained from square-integrable coefficients: $\mathcal{F}_2 = \left\{ f = Ta, ~~a \in L^2(\mathrm{d}\tau) \right\}$.
This can be realized by random Fourier features model \citep{rahimi2007random}.
which realizes the formulation of two-layer networks $f(\bm x) = \frac{1}{m} \sum_{i=1}^m a_i \phi(\bm x, \bm w_i)$.
If we only train the first layer, 
\fi

This contrast becomes explicit in random feature models (RFMs) \citep{rahimi2007random} and two-layer NNs, both of which can be written in the form the form $f(\bm x) = \frac{1}{m} \sum_{i=1}^m a_i \phi(\bm x, \bm w_i)$ with a feature map $\phi: \mathcal{X} \times \mathcal{W} \rightarrow \mathbb{R}$. The key difference lies in whether the features are learned: in RFMs, the weights $\{ \bm w_i\}_{i=1}^m \in \mathcal{W}$ are sampled i.i.d. from a probability measure $\mu$ and kept fixed, and only the second layer weights $\bm a :=\{ a_i \}_{i=1}^m$ are optimized, whereas in NNs both layers are jointly trained.

From a function-space perspective, RFMs with $\ell_2$-regularization on $\bm a$ are equivalent to kernel methods via the representer theorem \citep{scholkopf2001generalized}, with empirical kernel $\hat{k}(\bm x, \bm x') = \frac{1}{m} \sum_{i=1}^m \phi(\bm x, \bm w_i)\phi(\bm x', \bm w_i)$,
which approximates, as $m \to \infty$, the population kernel
\[
k_0(\bm x, \bm x') = \int_{\mathcal{W}} \phi(\bm x, \bm w)\phi(\bm x', \bm w)\, \mathrm{d}\mu(\bm w)\,.
\]
The associated reproducing kernel Hilbert space (RKHS), denoted $\mathcal{H}_0(\mu)$, is the closure of functions expressible as weighted averages of the features $\phi(\cdot, \bm w)$.

When the optimization of $\bm a$ is not $\ell_2$-regularized, the induced function space generally differs from an RKHS. For instance, under $\ell_p$-regularization with $1 \leq p < 2$ \citep{celentano2021minimum,chen2023duality}, the function space strictly enlarges relative to $\mathcal{H}_0(\mu)$ by Hölder’s inequality. In the limiting case of $\ell_1$-regularization, one recovers an $\mathcal{F}_1$-type space \citep{bach2017breaking}, closely related to the Barron space $\mathcal{B}$ \citep{barron1993universal,weinan2021barron}. The Barron space can be viewed as the largest function class that two-layer neural networks can learn efficiently in a statistical sense. The relationship between these spaces is clarified by \cite{chen2023duality}, which shows that $\mathcal{B} = \bigcup_{\mu \in \mathcal{P}(\mathcal{W})} \mathcal{H}_0(\mu)$, where $\mathcal{P}(\mathcal{W})$ denotes the set of probability measures on $\mathcal{W}$, establishing a natural connection between RFMs and two-layer neural networks. Notably, under the mean-field regime \citep{chizat2021convergence}, the function space induced by two-layer neural networks forms a subset of the Barron space \citep{wojtowytsch2020can}.

Despite existing characterizations of kernels and neural networks in function space, these perspectives are largely static: they do not explain how training algorithms (e.g., gradient descent) reshape the function space $\mathcal{H}_0(\mu)$ during feature learning. This gap motivates the following question:
\begin{center}
    \emph{How does the kernel (or function space) evolve under gradient updates, and what information is progressively learned?}
\end{center}
In this work, we address this question by providing a precise characterization of the function-space evolution induced by a single large gradient descent step in a two-layer network trained on a Gaussian single-index model. This setting, which has been extensively studied in recent theoretical work on feature learning, has been shown to capture the most general class of functions learnable in the proportional high-dimensional regime \citep{damian2022neural,ba2022high,dandi2023two}. Our analysis reveals how feature learning reshapes the underlying function space and improves its alignment with the target signal. Our contributions are as follows:

\paragraph{1. Approximation by a target-dependent Gaussian distribution:} In \cref{sec:three}, we prove that the post-update feature distribution is well approximated by a target-dependent spiked Gaussian, leading to a data-adaptive kernel $k_1$ (\emph{c.f.} \cref{thm:approx-true-kernel}). It demonstrates that feature learning can be expressed as target-dependent distribution transformation either in the parameter space or in the input-data space on such kernel.

\paragraph{2. Expansion of the data-adaptive kernel around an isotropic kernel:} In \cref{sec:four}, we study the spectrum of the kernel $k_1$. We conduct a Taylor expansion of $k_1$ around an isotropic kernel that isolates the contribution of the spike, and prove that higher-order remainder terms vanish as the dimension grows (\emph{c.f.} \cref{thm:cov-expansion}). In particular, these higher order terms take the form of isotropic kernels coupled with linear and non-linear projections of the input onto the target vector $\bm w^*$. This demonstrates that the role of feature learning is to impose an additional target-dependent kernels, thereby reshaping the function space. 

\paragraph{3. Feature learning in the top eigenspaces for the ReLU activation:} In \cref{sec:relu}, we consider the case of the ReLU activation function, to give an explicit characterization of the kernel $k_1$, its spectrum and the spanned function space. Our results theoretically prove that the spike transforms primarily the top and linear eigenspaces of the operator (\emph{c.f.} \cref{thm:lin-eigenfn},~\ref{thm:approx-top-eigenfn}). Our numerical validations support our theoretical findings and also illustrate the connection between data-adaptive kernels and neural networks.\\

% \begin{itemize}[leftmargin=*,topsep=0.5mm, itemsep=0.mm]
%     \item In \cref{sec:three}, we prove that the post-update feature distribution is well approximated by a target-dependent spiked Gaussian, leading to a data-adaptive kernel $k_1$ (\emph{c.f.} \cref{thm:approx-true-kernel}). It demonstrates that feature learning can be expressed as target-dependent distribution transformation either in the parameter space or in the input-data space on such kernel.
%     \item In \cref{sec:four}, we study the spectrum of the kernel $k_1$. We conduct a Taylor expansion of $k_1$ around an isotropic kernel that isolates the contribution of the spike, and prove that higher-order remainder terms vanish as the dimension grows (\emph{c.f.} \cref{thm:cov-expansion}). It shows that the role of feature learning is to impose an additional target-dependent kernel, thereby reshaping the function space. 
%     \item In \cref{sec:relu}, we consider the case of the ReLU activation function, to give an explicit characterization of the kernel, its spectrum and the spanned function space. Our results theoretically prove that the spike transforms primarily the top and linear eigenspaces of the operator (\emph{c.f.} \cref{thm:lin-eigenfn},~\ref{thm:approx-top-eigenfn}). Our numerical validations support our theoretical findings and also illustrate the connection between data-adaptive kernels and neural networks.
% \end{itemize}

These results provide a precise characterization of how the function space evolves during early training. In particular, they show that this evolution is modulated by the choice of step size, with larger step sizes inducing stronger transformations and increasing the contribution of both linear and higher-order non-linear features. This also bridges between data-adaptive kernels and neural networks, indicating a possibility of exploring proper initialization for feature learning.


\subsection{Related works}

The study of function spaces in deep learning theory stems from RKHS \citep{jacot2018neural_mini} and \citep{lee2017deep_mini}. However, these naturally operate in a lazy training regime \citep{chizat2019lazy_mini}, hence their ability for active feature learning is limited \citep{ghorbani2019limitations_mini}.

Recent works on feature learning in two-layer NNs are exemplified by how parameters change at early stages of training \citep{ba2022high,dandi2023two}. Collectively, these show that following a single step of gradient descent, the updated weight matrix $\bm W_1$ can be approximately decomposed into a deterministic rank-one signal component and a vanishing noise term. Moreover, \cite{moniri2023theory,cui2024asymptotics,dandi2024random} show this low-rank deformation injects informative spikes into the spectrum of the feature matrix $\sigma(\bm W_1^\top \bm x)$, contributing to the alignment towards the target function.
In parallel, \cite{Xu2024NeuralFL} studies feature learning via the \emph{feature geometry} perspective, unifying statistical dependence and feature representations in an inner-product space. However, their focus is on learning optimal features with standard networks rather than understanding how features evolve during training. Furthermore, \cite{dou2021training} studied the conjugate kernel RKHS $\mathcal{H}_t$ and the neural tangent kernel RKHS $\mathcal{K}_t$ through a signed measure. Nevertheless, their framework does not explicitly describe how the function space evolves after certain steps of gradient descent, nor which features are learned through optimization.

